\section{Segregate Odd and Even Position Nodes}
\label{sec:ll-odd-even}

\problemstatement{Given a linked list, group all nodes at odd positions
together followed by the nodes at even positions, preserving the relative
order within each group. Positions are 1-indexed; do it in place.}

\begin{patternbox}
Rearranging by position parity is a \textbf{two-chain weave}: maintain an
odd chain and an even chain, appending alternate nodes to each, then splice
the even chain onto the tail of the odd chain. Only pointers move.
\end{patternbox}

\subsection*{Build two chains, then join}

Keep an \texttt{odd} pointer starting at the head and an \texttt{even}
pointer starting at the second node, plus a saved \texttt{evenHead}. Repeatedly
link \texttt{odd->next = even->next} (skip the even node) and advance
\texttt{odd}; then \texttt{even->next = odd->next} and advance
\texttt{even}. This splits the list into odd-position and even-position
chains in a single pass. Finally attach \texttt{evenHead} after the last odd
node.

\begin{ideabox}[Interleave by Skipping Every Other Node]
Advancing each chain by two positions (skipping the other parity) threads
the odd and even nodes apart without moving data. The saved
\texttt{evenHead} lets the two chains be reconnected at the end.
\end{ideabox}

\subsection*{Algorithm}

\begin{enumerate}
  \item If the list has fewer than $3$ nodes, return it unchanged.
  \item \texttt{odd = head}, \texttt{even = head->next},
        \texttt{evenHead = even}.
  \item While \texttt{even} and \texttt{even->next}: relink odd then even by
        skipping one node each; advance both.
  \item \texttt{odd->next = evenHead}; return the head.
\end{enumerate}

\subsection*{Dry run}

\dryrun{$1\to2\to3\to4\to5$.}

\begin{itemize}
  \item Odd chain $1\to3\to5$; even chain $2\to4$.
  \item Join: $1\to3\to5\to2\to4$.
\end{itemize}

\subsection*{C++ implementation}

\begin{lstlisting}
#include <bits/stdc++.h>
using namespace std;

struct ListNode {
    int val; ListNode* next;
    ListNode(int v) : val(v), next(nullptr) {}
};

ListNode* build(const vector<int>& a) {
    ListNode dummy(0), *t = &dummy;
    for (int x : a) { t->next = new ListNode(x); t = t->next; }
    return dummy.next;
}

ListNode* oddEvenList(ListNode* head) {
    if (!head || !head->next) return head;
    ListNode* odd = head;
    ListNode* even = head->next;
    ListNode* evenHead = even;
    while (even && even->next) {
        odd->next = even->next;   odd = odd->next;    // splice next odd
        even->next = odd->next;   even = even->next;  // splice next even
    }
    odd->next = evenHead;                            // attach evens after odds
    return head;
}

void print(ListNode* h) { for (; h; h = h->next) cout << h->val << " "; cout << "\n"; }

int main() {
    print(oddEvenList(build({1, 2, 3, 4, 5})));    // 1 3 5 2 4
    return 0;
}
\end{lstlisting}

\begin{complexitybox}
\textbf{Time Complexity.} $O(n)$ -- one pass. \textbf{Space Complexity.}
$O(1)$.
\end{complexitybox}

\begin{takeawaybox}
Grouping by position parity is a two-chain weave joined at the end -- a
pure pointer rearrangement. The ``split into interleaved chains, then
reconnect'' idea also underlies reordering and unzip-style list problems.
\end{takeawaybox}
